\section{Conclusion and Future Work}\label{sec:conclusion}

High-order finite and spectral element discretizations are widely used in scientific and engineering domains to support high-fidelity simulations. The resulting datasets combine local tensor-product structure with globally unstructured element connectivity, posing distinct challenges for efficient storage and data movement. Existing error-bounded lossy compressors are not designed for this data organization and often struggle to jointly achieve high compression ratios, strict error control, and high throughput.

To address these challenges, we develop two complementary compressors, FDWZ and MRCZ, within a unified element-local framework with strict pointwise error guarantees. FDWZ extends the SZ framework with distance-weighted multilevel interpolation and multi-tree Huffman encoding for nonuniform nodal layouts. MRCZ represents each element using quantized transform coefficients, corrects the resulting prediction with error-bounded residuals, and reorganizes coefficients by basis mode to improve lossless coding. Both methods employ parallel element-centric processing while avoiding costly cross-element dependencies.

Experiments on large-scale high-order element datasets from multiple simulation software packages demonstrate substantial benefits for both storage reduction and end-to-end data movement. At matched PSNRs in the large-scale experiments, FDWZ improves the compression ratio over the best baseline by 80.8\% on average, and MRCZ further improves it over FDWZ by 111.4\%. Across the throughput experiments, FDWZ ranks among the two fastest methods in 38 of 40 cases. In the cross-facility data-movement experiments, FDWZ matches or outperforms the fastest baseline, reducing end-to-end time by up to 19.2\%. MRCZ retains competitive end-to-end time while producing the smallest transfers. At matched compression ratios, MRCZ and FDWZ also exhibit the smallest reconstruction errors in the visual comparison.

Future work will investigate a broader range of invertible and numerically stable transforms beyond the Chebyshev-based transform adopted in this study. We also plan to develop GPU implementations of FDWZ and MRCZ to accelerate their prediction, quantization, transform, and lossless-encoding stages, thereby improving throughput and scalability on modern heterogeneous computing platforms.
